\newpage
\section{the initial wave packet of the projectile}\label{sec:proj-wavepacket}

the results of this chapter depend only on the free particle (the ``projectile'', coordinate $X$, mass $M$) and not on the internal structure of the atom(s) it scatters from. They are therefore used unchanged in Part~\ref{part:continuous-atoms} (continuous-position atoms) and in Part~\ref{part:discrete-atoms} (discrete-position atoms).

the projectile is prepared, at $t=t_i$, in the symmetric bimodal wave packet
\begin{align}
  \psi_\alpha(X) = \frac{1}{\sqrt{\sigma_\alpha\sqrt{\pi}}}\sqrt{\frac{2}{1+e^{-\frac{P_\alpha^2\sigma_\alpha^2}{\hbar^2}}}}e^{-\frac{X^2}{2\sigma_\alpha^2}}\cos\left(\frac{P_\alpha X}{\hbar}\right)\label{eq:proj-psi-alpha}
\end{align}
i.e.\ a superposition of two Gaussian packets of width $\sigma_\alpha$ with mean momenta $\pm P_\alpha$, normalized so that $\int|\psi_\alpha(X)|^2dX=1$ (see Sect.~\ref{sec:app-gaussian-normalization} and Sect.~\ref{sec:app-bimodal-normalization} for the normalization details).

the energy eigenfunction (momentum) expansion of $\psi_\alpha$ is
\begin{align}
  \psi_\alpha(X)=\frac1{\sqrt{2\pi\hbar}}\int\limits_{-\infty}^{\infty}dP\, C(P)e^{\frac{iPX}{\hbar}}
\end{align}
where
\begin{align}
  C(P) = \langle P|\psi_\alpha\rangle = \frac{1}{\sqrt{2\pi\hbar}}\int\limits_{-\infty}^{\infty}dX\, e^{-\frac{i}{\hbar}PX}\psi_\alpha(X)
\end{align}
Substituting the explicit form of $\psi_\alpha(X)$ yields:
\begin{align}
  C(P) = \frac{1}{\sqrt{2\pi\hbar}}\frac{1}{\sqrt{\sigma_\alpha\sqrt{\pi}}}\sqrt{\frac{2}{1+e^{-\frac{P_\alpha^2\sigma_\alpha^2}{\hbar^2}}}}\int\limits_{-\infty}^{\infty}dX\, e^{-\frac{i}{\hbar}PX}e^{-\frac{X^2}{2\sigma_\alpha^2}}\cos\left(\frac{P_\alpha X}{\hbar}\right)
\end{align}
Using $\cos(u) = \frac{1}{2}(e^{iu} + e^{-iu})$, we can evaluate the integral:
\begin{align}
  \int\limits_{-\infty}^{\infty}dX\, e^{-\frac{i}{\hbar}PX}e^{-\frac{X^2}{2\sigma_\alpha^2}}\cos\left(\frac{P_\alpha X}{\hbar}\right) & = \frac{1}{2}\int\limits_{-\infty}^{\infty}dX\, e^{-\frac{X^2}{2\sigma_\alpha^2}}\left[ e^{-\frac{i}{\hbar}(P-P_\alpha)X} + e^{-\frac{i}{\hbar}(P+P_\alpha)X} \right] \nonumber \\
  & = \frac{1}{2}\sqrt{2\pi\sigma_\alpha^2}\left[ e^{-\frac{(P-P_\alpha)^2\sigma_\alpha^2}{2\hbar^2}} + e^{-\frac{(P+P_\alpha)^2\sigma_\alpha^2}{2\hbar^2}} \right]
\end{align}
Substituting this back into the expression for $C(P)$ and simplifying the coefficients, we obtain the final result:
\begin{importantbox}
  \begin{align}
    C(P) = \sqrt{\frac{\sigma_\alpha}{2\hbar\sqrt{\pi}\left(1+e^{-\frac{P_\alpha^2\sigma_\alpha^2}{\hbar^2}}\right)}}\left[ e^{-\frac{(P-P_\alpha)^2\sigma_\alpha^2}{2\hbar^2}} + e^{-\frac{(P+P_\alpha)^2\sigma_\alpha^2}{2\hbar^2}} \right]\label{eq:proj-CP}
  \end{align}

  Note that for $P_\alpha \gg \hbar/\sigma_\alpha$, the second term in the brackets is negligible, and we can approximate $C(P)$ as:
  \begin{align}
    C(P) \approx \sqrt{\frac{\sigma_\alpha}{2\hbar\sqrt{\pi}}}\left[ e^{-\frac{(P-P_\alpha)^2\sigma_\alpha^2}{2\hbar^2}}+ e^{-\frac{(P+P_\alpha)^2\sigma_\alpha^2}{2\hbar^2}} \right]\label{eq:proj-CP-approx}
  \end{align}
\end{importantbox}
